\begin{proof}

\begin{enumerate}[label={$\langle1\rangle$\arabic*.},leftmargin=*]
\item There are finitely many critical values: the critical locus is analytic, its image is analytic by properness, and it does not contain a general point of the curve. Away from these values $g$ is a differentiable fibre bundle whose fibre has Euler number zero. Euler characteristic is additive over this open base and the finite exceptional set, so $e(Y)=\sum_c e((Y_c)_{\rm red})$. Nilpotents leave the underlying space unchanged. An irreducible support has a single multiplicity, so \ref{cdp:st:C.2.NC.3} applies; a reducible support is covered by \ref{cdp:st:C.2.NC.4}. A smooth support has Euler number zero directly from (C.NC.1). Thus every summand equals its triple-point count. \ref{cdp:st:C.2.NC.4} identifies the positive summands and their value two. Gauss--Bonnet gives $c_3(Y)=e(Y)$. Poincaré duality gives $e(Y)=2-2b_1(Y)+2b_2(Y)-b_3(Y)$, proving the assertion about $r$ and $b_3$. \ref{cdp:st:C.3.5} proves the required zero Euler number for every smooth fibre under the CDP hypotheses, so that application uses no assumption that the general fibre is a torus.

\item For a positive support, every line bundle on $Y$ pulls back trivially to its smooth rational normalization: its rational first Chern class is zero, and the injectivity and torsion-freeness used in \ref{cdp:st:C.2.NC.4} apply. In particular $\eta^*((K_Y^{-1}\otimes L)|_S)$ is trivial. Serre duality on the connected smooth compact surface $T$ identifies the first group in \eqref{cdp:eq:C.NC.11} with $H^0(T,\OO_T)^\vee\simeq\C$.

\item For the ambient group set $A=(K_Y\otimes L^{-1})|_S$ and $N^*=\OO_Y(-S)|_S$. Both $\eta^*A$ and $\eta^*N^*$ are trivial, even when the full fibre is $mS$. The injection of torsion-free differentials into normalization differentials gives
$H^0(S,\widetilde\Omega_S^1\otimes A)\hookrightarrow H^0(T,\Omega_T^1)=0$.
The canonical normal-crossings kernel isomorphism \ref{cdp:st:A.1.2} therefore gives
\[
 H^0(S,\Omega_Y^1|_S\otimes A)
 =H^0(T,\eta^*(N^*\otimes A))=\C.
\]
This proves the second equality. It establishes the universal normalization target requested in B.3 for the full normal-crossings class, including unequal multiplicities. It leaves the unrestricted nonordinary-fibre assertion to be decided separately.

\end{enumerate}
\end{proof}
